\documentclass[a4paper]{article}
\usepackage[fontset=fandol]{ctex}
\usepackage{amsmath, amssymb}
\usepackage{graphicx}
\usepackage[margin=2.5cm]{geometry}
\usepackage[most]{tcolorbox}
\usepackage{etoolbox}
\usepackage{listings}
\usepackage{hyperref}
\usepackage{booktabs}
\usepackage{subcaption}
\usepackage{float}
\usepackage{tikz}
\IfFileExists{pgfplots.sty}{\usepackage{pgfplots}\pgfplotsset{compat=1.18}}{}
\newtcolorbox{knowledgebox}[1]{enhanced, colback=blue!5!white, colframe=blue!75!black, colbacktitle=blue!75!black, coltitle=white, fonttitle=\bfseries, title=#1, attach boxed title to top left={yshift=-2mm, xshift=2mm}, boxrule=1pt, sharp corners}
\newtcolorbox{importantbox}[1]{enhanced, colback=yellow!10!white, colframe=yellow!80!black, colbacktitle=yellow!80!black, coltitle=black, fonttitle=\bfseries, title=#1, sharp corners}
\newtcolorbox{warningbox}[1]{enhanced, colback=red!5!white, colframe=red!75!black, colbacktitle=red!75!black, coltitle=white, fonttitle=\bfseries, title=#1, sharp corners}
\lstset{language=Python,basicstyle=\ttfamily\small,keywordstyle=\color{blue},stringstyle=\color{red!60!black},commentstyle=\color{green!60!black},breaklines=true,frame=single,numbers=left,numberstyle=\tiny\color{gray},captionpos=b,extendedchars=false}
\newcommand{\notetitle}{Advanced Topics on End-to-End ASR II}
\newcommand{\noteauthors}{Codex 基于课程页、公开视频字幕、coverage ledger 与补充官方材料整理}
\newcommand{\notedate}{\today}
\newcommand{\videochannel}{WAVLab}
\newcommand{\videopublishdate}{2023-11-20}
\newcommand{\videoduration}{unknown}
\newcommand{\videourl}{https://www.youtube.com/watch?v=MtwolOGs0\_A}
\newcommand{\videocoverpath}{cover.jpg}
\begin{document}
\begin{titlepage}
\centering
{\Large 课程笔记\par}
\vspace{1.2cm}
{\huge\bfseries \notetitle\par}
\vspace{0.8cm}
{\large \noteauthors\par}
\vspace{0.3cm}
{\large \notedate\par}
\vspace{1.2cm}
\ifdefempty{\videocoverpath}{{\small 公开视频封面不可用。\par}}{\includegraphics[width=0.82\textwidth,height=0.45\textheight,keepaspectratio]{\videocoverpath}\par}
\vfill
\begin{tcolorbox}[width=0.9\textwidth, colback=black!2!white, colframe=black!60, sharp corners]
\textbf{视频作者/频道}：\videochannel\par
\textbf{发布日期}：\videopublishdate\par
\textbf{视频时长}：\videoduration\par
\textbf{视频链接}：\href{\videourl}{\nolinkurl{\videourl}}
\end{tcolorbox}
\end{titlepage}
\tableofcontents
\newpage
\section{Multilingual ASR recap and constraints}
根据课程页主题、公开视频字幕以及本讲的 coverage ledger，Multilingual ASR recap and constraints 这一部分主要围绕 Advanced Topics on End-to-End ASR II | Advanced topics on end-to-end ASR II；multilingual ASR 的继续讨论：为什么字典 / 词表假设在很多语言里并不成立。 展开。
写作上保持 coverage-first 原则：先把老师在这一阶段真正讲过的问题设定、模型部件、训练或推理逻辑讲清，再压缩成适合复习的结构。

\begin{figure}[H]
\centering
\includegraphics[width=0.92\textwidth]{figures/coverage_map.png}
\caption{本讲 coverage map：按 coverage\_units.jsonl 归纳出的核心主题与章节映射。}
\end{figure}

\begin{figure}[H]
\centering
\includegraphics[width=0.92\textwidth]{figures/timeline.png}
\caption{本讲 timeline：根据字幕时间范围归纳的主题推进顺序。}
\end{figure}

\subsection{Advanced Topics on End-to-End ASR II}
从证据位置看，这一部分对应的课程证据区间是 2023-11-20。课程在这里强调的核心内容可以概括为：Advanced Topics on End-to-End ASR II | Advanced topics on end-to-end ASR II。
从课程组织上看，这一部分承担的是把局部概念嵌回整条识别流水线中的作用，避免读者把公式或模块孤立地记忆。

在本章结构中，这一单元被并入“Multilingual ASR recap and constraints”，目的是把同一阶段的定义、机制和权衡集中讲清。

\subsection{multilingual ASR 的继续讨论}
从证据位置看，这一部分对应的课程证据区间是 00:00:00,840--00:07:21,080。课程在这里强调的核心内容可以概括为：multilingual ASR 的继续讨论：为什么字典 / 词表假设在很多语言里并不成立。。
这一段更像课程的现代延伸：它把传统 ASR 课程内容与 multilingual transfer、speech foundation model、fine-tuning 实践接到了一起。

在本章结构中，这一单元被并入“Multilingual ASR recap and constraints”，目的是把同一阶段的定义、机制和权衡集中讲清。

\subsection{本章小结}
Multilingual ASR recap and constraints 这一部分把 Advanced Topics on End-to-End ASR II 中相邻的教学单元串成了一条连续叙事：既交代了这一阶段要解决的问题，也交代了方法、推理或工程权衡为什么要这样组织。

\section{Whisper-style speech foundation models}
根据课程页主题、公开视频字幕以及本讲的 coverage ledger，Whisper-style speech foundation models 这一部分主要围绕 Whisper-style speech foundation model 的开源设计、模型边界与主要取舍。 展开。
写作上保持 coverage-first 原则：先把老师在这一阶段真正讲过的问题设定、模型部件、训练或推理逻辑讲清，再压缩成适合复习的结构。

\subsection{Whisper-style speech foundation model 的开源设计、模型边界与主要取舍}
从证据位置看，这一部分对应的课程证据区间是 00:07:18,120--00:15:07,120。课程在这里强调的核心内容可以概括为：Whisper-style speech foundation model 的开源设计、模型边界与主要取舍。。
这一段更像课程的现代延伸：它把传统 ASR 课程内容与 multilingual transfer、speech foundation model、fine-tuning 实践接到了一起。

记录系统中的附加说明指出：CS224S supplement anchor: foundation models for spoken language and speech recognition beyond English.

在本章结构中，这一单元被并入“Whisper-style speech foundation models”，目的是把同一阶段的定义、机制和权衡集中讲清。
\begin{knowledgebox}{CS224S 补充}
CS224S supplement anchor: foundation models for spoken language and speech recognition beyond English.
\end{knowledgebox}

\subsection{本章小结}
Whisper-style speech foundation models 这一部分把 Advanced Topics on End-to-End ASR II 中相邻的教学单元串成了一条连续叙事：既交代了这一阶段要解决的问题，也交代了方法、推理或工程权衡为什么要这样组织。

\section{Inference details and timestamps}
根据课程页主题、公开视频字幕以及本讲的 coverage ledger，Inference details and timestamps 这一部分主要围绕 模型 architecture 的修改、timestamp / window shifting 等推理细节。 展开。
写作上保持 coverage-first 原则：先把老师在这一阶段真正讲过的问题设定、模型部件、训练或推理逻辑讲清，再压缩成适合复习的结构。

\subsection{模型 architecture 的修改、timestamp / window shifting 等推理细节}
从证据位置看，这一部分对应的课程证据区间是 00:15:04,600--00:21:57,559。课程在这里强调的核心内容可以概括为：模型 architecture 的修改、timestamp / window shifting 等推理细节。。
从课程组织上看，这一部分承担的是把局部概念嵌回整条识别流水线中的作用，避免读者把公式或模块孤立地记忆。

记录系统中的附加说明指出：CS224S supplement anchor: compare with modern ASR encoders, Whisper inference, and multilingual decoding strategies.

在本章结构中，这一单元被并入“Inference details and timestamps”，目的是把同一阶段的定义、机制和权衡集中讲清。
\begin{knowledgebox}{CS224S 补充}
CS224S supplement anchor: compare with modern ASR encoders, Whisper inference, and multilingual decoding strategies.
\end{knowledgebox}

\subsection{本章小结}
Inference details and timestamps 这一部分把 Advanced Topics on End-to-End ASR II 中相邻的教学单元串成了一条连续叙事：既交代了这一阶段要解决的问题，也交代了方法、推理或工程权衡为什么要这样组织。

\section{Open-source design choices and results}
根据课程页主题、公开视频字幕以及本讲的 coverage ledger，Open-source design choices and results 这一部分主要围绕 开源 demo、实验结果与具体 design choices，说明大模型在多语种场景中的实际行为。 展开。
写作上保持 coverage-first 原则：先把老师在这一阶段真正讲过的问题设定、模型部件、训练或推理逻辑讲清，再压缩成适合复习的结构。

\subsection{开源 demo、实验结果与具体 design choices，说明大模型在多语种场景中的实际行为}
从证据位置看，这一部分对应的课程证据区间是 00:21:54,520--00:28:46,000。课程在这里强调的核心内容可以概括为：开源 demo、实验结果与具体 design choices，说明大模型在多语种场景中的实际行为。。
从课程组织上看，这一部分承担的是把局部概念嵌回整条识别流水线中的作用，避免读者把公式或模块孤立地记忆。

在本章结构中，这一单元被并入“Open-source design choices and results”，目的是把同一阶段的定义、机制和权衡集中讲清。

\subsection{本章小结}
Open-source design choices and results 这一部分把 Advanced Topics on End-to-End ASR II 中相邻的教学单元串成了一条连续叙事：既交代了这一阶段要解决的问题，也交代了方法、推理或工程权衡为什么要这样组织。

\section{Scaling and open questions}
根据课程页主题、公开视频字幕以及本讲的 coverage ledger，Scaling and open questions 这一部分主要围绕 scaling、emergent capabilities 与 speech 领域的开放问题。 展开。
写作上保持 coverage-first 原则：先把老师在这一阶段真正讲过的问题设定、模型部件、训练或推理逻辑讲清，再压缩成适合复习的结构。

\subsection{scaling、emergent capabilities 与 speech 领域的开放问题}
从证据位置看，这一部分对应的课程证据区间是 00:28:42,919--00:35:51,040。课程在这里强调的核心内容可以概括为：scaling、emergent capabilities 与 speech 领域的开放问题。。
从课程组织上看，这一部分承担的是把局部概念嵌回整条识别流水线中的作用，避免读者把公式或模块孤立地记忆。

记录系统中的附加说明指出：CS224S supplement anchor: foundation models, multilingual ASR, and low-resource language varieties.

在本章结构中，这一单元被并入“Scaling and open questions”，目的是把同一阶段的定义、机制和权衡集中讲清。
\begin{knowledgebox}{CS224S 补充}
CS224S supplement anchor: foundation models, multilingual ASR, and low-resource language varieties.
\end{knowledgebox}

\subsection{本章小结}
Scaling and open questions 这一部分把 Advanced Topics on End-to-End ASR II 中相邻的教学单元串成了一条连续叙事：既交代了这一阶段要解决的问题，也交代了方法、推理或工程权衡为什么要这样组织。

\section{总结与延伸}
Advanced Topics on End-to-End ASR II 这一讲在整门《Speech Recognition and Understanding》课程中的位置，是把 Advanced E2E ASR II 讲成一个可复用的建模模块。
从教材化角度看，本讲最重要的不是记住单一术语，而是理解它如何嵌入完整的 automatic speech recognition pipeline：输入语音如何被表示、模型如何被训练、解码如何得到最终转写，以及这些设计分别带来什么假设和权衡。
由于公开源以课程页和公开视频字幕为主、缺少官方 slide PDF，本章在正文里尽量把术语、方法和工程语境解释完整，并把缺失项留在 omission 和 manifest 里，而不是凭空补写。
\end{document}
